\documentclass[12pt]{article}
\usepackage[utf8]{inputenc}
\usepackage[T1]{fontenc}
\usepackage{hyperref}
\title{Uradhi Language Resources for K-12 Curriculum Development}
\author{Compiled by Claude Sonnet 5 (claude-sonnet-5) for Andr\'as Kornai}
\date{}
\begin{document}
\maketitle
\section{Introduction}
Uradhi (ISO 639-3 urf; Glottocode urad1239) is an Australian language of northern Cape York, Queensland, with the Atampaya dialect. Sources conflict on its status: Wikipedia lists it as apparently extinct (extinction date 1990s, revival 2015), while another source says a few partial speakers of Atampaya remain. Local speakers call their language Injinoo Ikya, a term shared with Angkamuthi and Yadhaykenu. Little classroom material was found.
\section{Primary Educational Resources}
None found.
\section{Linguistic and Reference Materials}
\begin{itemize}
\item Uradhi language, Wikipedia: \url{https://en.wikipedia.org/wiki/Uradhi_language}
\item Uradhi dialect cluster, About the North: \url{https://aboutthenorth.au/styled-59/styled-184/styled-238/styled-472/}
\item Hale (1976) reference record, Glottolog: \url{https://glottolog.org/resource/reference/id/96843}
\end{itemize}
\section{Religious and Community Resources}
\begin{itemize}
\item Pama Language Centre, Yadhaykenu page (a related variety, not Uradhi itself): \url{https://www.pamacentre.org.au/yadhaykenyu/}
\end{itemize}
\section{Digital and Audio Resources}
None found.
\section{Curriculum Development Context}
The Pama Language Centre is the likely regional partner for any program, and Injinoo Ikya community work would be the natural frame. The Pama page on Injinoo Ikya was listed but not read, so its content is unverified.
\end{document}